\documentclass[aspectratio=169,11pt]{beamer}
% Shared CMPSC 480 Beamer preamble.
\usepackage[T1]{fontenc}
\usepackage[utf8]{inputenc}
\usepackage{lmodern}
\usepackage{amsmath,amssymb,mathtools}
\usepackage{booktabs,array,tabularx}
\usepackage{xcolor}
\usepackage{listings}
\usepackage{tikz}
\usetikzlibrary{arrows.meta,positioning,shapes.geometric}

% Ignore only negligible TeX box diagnostics that are below visual tolerance.
% Larger layout defects still appear in the build log and in rendered QA.
\vfuzz=3pt
\hbadness=2000

\definecolor{courseink}{HTML}{111111}
\definecolor{coursegray}{HTML}{EDEDED}
\definecolor{courserule}{HTML}{B8BCC4}
\definecolor{courseblue}{HTML}{3D8DFF}
\definecolor{courselightblue}{HTML}{D0EDFA}
\definecolor{coursemuted}{HTML}{5C6470}
\definecolor{coursegreen}{HTML}{0B7A53}
\definecolor{coursered}{HTML}{B3261E}

\usetheme{default}
\usefonttheme{professionalfonts}
\setbeamertemplate{navigation symbols}{}
\setbeamersize{text margin left=0.65cm,text margin right=0.65cm}

\setbeamercolor{normal text}{fg=courseink,bg=white}
\setbeamercolor{frametitle}{fg=courseink,bg=white}
\setbeamercolor{title}{fg=courseink,bg=white}
\setbeamercolor{subtitle}{fg=coursemuted,bg=white}
\setbeamercolor{structure}{fg=courseblue}
\setbeamercolor{alerted text}{fg=coursered}
\setbeamercolor{block title}{fg=courseink,bg=courselightblue}
\setbeamercolor{block body}{fg=courseink,bg=coursegray}
\setbeamercolor{example text}{fg=coursegreen}

\setbeamerfont{title}{size=\fontsize{25}{29}\selectfont,series=\bfseries}
\setbeamerfont{subtitle}{size=\fontsize{13}{16}\selectfont}
\setbeamerfont{frametitle}{size=\fontsize{19}{22}\selectfont,series=\bfseries}
\setbeamerfont{normal text}{size=\fontsize{11}{14}\selectfont}
\setbeamerfont{footline}{size=\fontsize{7.5}{9}\selectfont}

\setbeamertemplate{frametitle}{%
  \nointerlineskip
  % Use content-driven height so a long assessment title can wrap without
  % being clipped by a fixed-height title bar.
  \begin{beamercolorbox}[wd=\paperwidth,sep=0.17cm,leftskip=0.48cm,rightskip=0.48cm]{frametitle}
    \insertframetitle
  \end{beamercolorbox}
  \vspace{-0.08cm}
  {\color{courserule}\rule{\textwidth}{0.35pt}}
}

\setbeamertemplate{footline}{%
  \leavevmode
  \begin{beamercolorbox}[wd=\paperwidth,ht=2.6ex,dp=1.1ex,leftskip=.65cm,rightskip=.65cm]{author in head/foot}
    \color{coursemuted}\insertshorttitle\hfill\insertframenumber/\inserttotalframenumber
  \end{beamercolorbox}
}

\setbeamertemplate{itemize item}{\color{courseblue}\small$\blacktriangleright$}
\setbeamertemplate{itemize subitem}{\color{courseblue}\scriptsize$\bullet$}
\setbeamertemplate{enumerate item}{\color{courseblue}\insertenumlabel.}

\lstdefinestyle{coursepython}{
  language=Python,
  basicstyle=\ttfamily\fontsize{8.5}{10}\selectfont,
  keywordstyle=\color{courseblue}\bfseries,
  commentstyle=\color{coursemuted}\itshape,
  stringstyle=\color{coursegreen},
  showstringspaces=false,
  columns=fullflexible,
  keepspaces=true,
  frame=single,
  rulecolor=\color{courserule},
  backgroundcolor=\color{coursegray},
  breaklines=true,
  tabsize=4,
  xleftmargin=0.15cm,
  xrightmargin=0.15cm,
  framexleftmargin=0.15cm,
  framexrightmargin=0.15cm
}
\lstset{style=coursepython}

\newcommand{\points}[1]{\hfill{\color{courseblue}\textbf{[#1 points]}}}
\newcommand{\deliverable}[1]{\textbf{\color{courseblue}#1}}
\newcommand{\code}[1]{\texttt{#1}}
\newcommand{\keyidea}[1]{%
  \begin{beamercolorbox}[rounded=false,sep=0.55em,wd=\linewidth]{block body}
    \textbf{Key idea.} #1
  \end{beamercolorbox}
}
\newcommand{\answerlabel}{\textbf{\color{coursegreen}Answer.}\ }
\newcommand{\gradinglabel}{\textbf{\color{courseblue}Grading guidance.}\ }

\newenvironment{questionblock}[1]
  {\begin{block}{#1}}
  {\end{block}}

\newenvironment{solutionblock}[1]
  {\begin{block}{#1}}
  {\end{block}}

\AtBeginSection[]{
  \begin{frame}
    \vfill
    {\usebeamerfont{title}\color{courseink}\insertsectionhead\par}
    \vspace{0.4cm}
    {\color{courseblue}\rule{0.32\paperwidth}{3pt}}
    \vfill
  \end{frame}
}


% CMPSC480_TOTAL_POINTS: 10
% CMPSC480_ITEM_IDS: 1,2,3,4
\title[Week 6 Assignment Solutions]{Week 6 Assignment: Linear Q-Function Approximation}
\subtitle{Feature design, semi-gradient updates, and stability evidence}
\author{CMPSC 480 - Instructor Solution Deck}
\date{}

\begin{document}


\begin{frame}
  \titlepage
\vspace{-0.35cm}
\begin{center}
\color{coursered}\bfseries Instructor material - do not distribute with the question deck
\end{center}
\end{frame}

\begin{frame}{Solution overview}
\begin{columns}[T,onlytextwidth]
\begin{column}{0.62\textwidth}
\Large This instructor deck provides a complete matched solution and grading guidance.

\vspace{0.35cm}
\normalsize
Coverage: Week 6: function approximation, semi-gradients, and the deadly triad
\end{column}
\begin{column}{0.33\textwidth}
\begin{block}{Points}10 total\end{block}
\begin{block}{Expected time}6-8 hours\end{block}
\begin{block}{Submission}Repository plus experiment memo\end{block}
\end{column}
\end{columns}
\end{frame}

\begin{frame}{Instructor verification checklist}
\small
\begin{itemize}
  \item Use Python 3.11 and NumPy; do not use a library that implements the assessed algorithm.
  \item Keep data, model, optimization, and evaluation responsibilities behind explicit interfaces.
  \item Validate shapes, ranges, finite values, empty inputs, and terminal or degenerate cases.
  \item Use deterministic seeds and record every configuration used for reported evidence.
  \item Include focused tests and a README with exact reproduction commands.
\end{itemize}
\end{frame}

\begin{frame}{Solution 1.a - Features and dimensions (1 points)}
\small
\textbf{Question focus.} Define at least four state-action features and justify their scales.

\vspace{0.25cm}
\answerlabel A full-credit design names each signal, such as normalized coordinates, goal distance, wall risk, and action alignment; all are finite and on comparable documented scales.
\end{frame}

\begin{frame}{Solution 1.b - Features and dimensions (1 points)}
\small
\textbf{Question focus.} Prove by an executable assertion that parameter size does not grow when the environment map grows.

\vspace{0.25cm}
\answerlabel Instantiate two map sizes with the same feature schema and assert equal feature and weight shapes. Reject any feature vector with a mismatched dimension.
\end{frame}

\begin{frame}{Solution 2.a - Linear values and policy (1 points)}
\small
\textbf{Question focus.} Compute Q-hat(s,a;w) and test a known vector by hand.

\vspace{0.25cm}
\answerlabel Return float(w dot f). For w=[1,-2,0.5] and f=[2,1,4], Q-hat=2-2+2=2.
\end{frame}

\begin{frame}{Solution 2.b - Linear values and policy (1 points)}
\small
\textbf{Question focus.} Specify training and evaluation tie behavior.

\vspace{0.25cm}
\answerlabel Training explores uniformly with probability epsilon and otherwise uses a stable argmax over legal actions. Evaluation sets epsilon to zero and uses the same documented stable tie order.
\end{frame}

\begin{frame}{Solution 3.a - Semi-gradient update (2 points)}
\small
\textbf{Question focus.} Derive the update from one-half times squared TD error.

\vspace{0.25cm}
\answerlabel With y treated as fixed, negative loss gradient is (y-Q-hat) grad Q-hat. For a linear function grad Q-hat=f, so w receives alpha times delta times f.
\end{frame}

\begin{frame}{Solution 3.b - Semi-gradient update (1 points)}
\small
\textbf{Question focus.} Compute an update for w=[0,0], f=[1,2], reward=1, gamma=0.5, max-next-Q=4, and alpha=0.1.

\vspace{0.25cm}
\answerlabel The target is 3, current Q is 0, delta is 3, and the new weights are [0.3,0.6].
\end{frame}

\begin{frame}{Solution 3.c - Semi-gradient update (1 points)}
\small
\textbf{Question focus.} State and test the terminal target.

\vspace{0.25cm}
\answerlabel When done is true, y equals reward with no next-state maximum. A test uses an enormous next Q and confirms the terminal update is unchanged.
\end{frame}

\begin{frame}{Solution 4.a - Experiment and stability (1 points)}
\small
\textbf{Question focus.} Run one feature ablation and compare with a tabular small-state oracle.

\vspace{0.25cm}
\answerlabel Hold seed, budget, and hyperparameters fixed; report return and action agreement. Explain whether the removed feature changes generalization or bias.
\end{frame}

\begin{frame}{Solution 4.b - Experiment and stability (1 points)}
\small
\textbf{Question focus.} Define an instability gate using at least three logged quantities.

\vspace{0.25cm}
\answerlabel A suitable gate watches nonfinite values, weight norm, TD-error quantiles, and evaluation return across seeds, and stops or flags the run when declared thresholds are exceeded.
\end{frame}

\begin{frame}{Edge-case reasoning and expected behavior}
\small
\begin{itemize}
  \item The prediction is zero and no coordinate changes.
  \item Target equals reward only.
  \item Raise a clear error before the dot product.
  \item Use stable evaluation order.
  \item Stop, record the transition and configuration, and mark the run invalid.
\end{itemize}
\end{frame}

\begin{frame}{Grading rubric}
\scriptsize
\begin{tabularx}{\textwidth}{>{\bfseries}p{0.28\textwidth} p{0.08\textwidth} X}
\toprule
Category & Points & Full-credit evidence \\
\midrule
Features and dimensions & 2 & Features are fixed length, normalized, documented, and safely validated. \\
Linear values and policy & 2 & Dot product, legal-action filtering, probabilities, and ties are correct. \\
Semi-gradient update & 4 & Target, TD error, gradient, terminal rule, and signs are correct. \\
Experiment and stability & 2 & Controlled comparison includes return, TD error, parameter norm, and multiple seeds. \\
\bottomrule
\end{tabularx}
\end{frame}

\begin{frame}{Reflection exemplar}
\small
\answerlabel A strong response links feature overlap to which predictions changed, cites controlled evidence, and distinguishes improved return from stable value estimates.

\vspace{0.25cm}
\gradinglabel Award credit for a specific claim supported by technical evidence from the submitted work.
\end{frame}

\end{document}
